Checked in August 2026. Venues close and events move, so search before you travel.

\subsection{The one that is coming up}\label{li02_scene-the-one-that-is-coming-up}

\textbf{ARC Music Festival}, Union Park, \textbf{4 to 7 September 2026}. Expanded to four days this
year. It is the big one, it is in the West Loop, and it deliberately programmes Chicago house
originators alongside international techno acts.

If you are going to do one thing this year to connect with the music, this is the
low-friction option, and it is about three weeks away as of writing.

\subsection{Clubs}\label{li02_scene-clubs}

\textbf{Smartbar}, 3730 N Clark St, underneath Metro in Lakeview. Opened 1982, capacity around
400, still independent, and the only club from the original house era still operating. Frankie
Knuckles, Derrick Carter and The Blessed Madonna have all held residencies. This is the
single most important room in the city for what you are trying to learn, and a Tuesday night
there costs very little.

\textbf{Spybar}, River North. Long-running, techno-leaning.

\textbf{Sound-Bar}, River North. Bigger, more mainstream.

\textbf{Primary}, West Town. House and techno, later hours.

\textbf{The Promontory}, Hyde Park. More live and soul-oriented, with house nights.

\textbf{Cermak Hall / Radius}, Pilsen. Large-format events.

A practical note: a lot of the best Chicago house happens at one-off parties and loft events
rather than at fixed venues, and those are advertised on Resident Advisor and Instagram
rather than anywhere findable. Following a few local DJs and promoters is more useful than
checking club calendars.

\subsection{Annual events}\label{li02_scene-annual-events}

\textbf{Chosen Few Picnic}, Jackson Park, usually the first or second Saturday of July. The 2026
edition was 11 July, so the next is summer 2027. Known as the Woodstock of house music, it
is free to attend in the general area, it has been running since 1990, and it is the closest
thing the genre has to a family reunion. The Chosen Few DJs include Wayne Williams, Jesse
Saunders, Terry Hunter and Alan King.

If you only go to one thing ever, go to this.

\textbf{Chicago House Music Festival and Conference}, run by the city, free admission, usually
late May. The conference component has panels with people who were actually there in 1985.

\textbf{Silver Room Block Party}, Hyde Park, summer.

\subsection{Record shops}\label{li02_scene-record-shops}

\textbf{Gramaphone Records}, 2843 N Clark St, Lakeview. Trading since 1969 and focused on dance
music since house took hold in the early 1980s. Deep 12-inch inventory, classics, imports and
current club records. The staff are DJs and will talk to you if you ask sensible questions.

\textbf{Dusty Groove}, Ukrainian Village. Broader, strong on soul, disco and jazz, which is where
house came from.

\textbf{Reckless Records}, several locations. General, but worth digging.

You do not need turntables to benefit from a record shop. Going in, asking what is good, and
writing down names is a legitimate use of the place.

\subsection{Radio and online}\label{li02_scene-radio-and-online}

\textbf{WHPK 88.5 FM}, University of Chicago student radio, with long-running dance and house
programming.

\textbf{Chicago House FM}, online.

\textbf{5 Magazine} (5mag.net) covers the Chicago scene specifically and has done for years. Their
event calendar is more reliable than anything else for local house.

\textbf{Resident Advisor} for club listings, and \textbf{Do312} for general Chicago events.

\subsection{Where to buy the music}\label{li02_scene-where-to-buy-the-music}

\textbf{Bandcamp} pays artists better than anything else and Chicago producers use it heavily.

\textbf{Beatport} for DJ-oriented releases, with the key and BPM listed on every track, which is
useful reference data even if you never DJ.

\textbf{Discogs} for finding out what a record actually is, who played on it, and what the pressing
history was. The database is extraordinary.

\subsection{If you want to meet people}\label{li02_scene-if-you-want-to-meet-people}

Two honest observations.

The Chicago house community is welcoming to people who are genuinely interested and
noticeably less welcoming to people who treat it as a genre to consume rather than a
tradition with living participants. Knowing the names in \texttt{chicago\_canon.md} goes a long way.

And the fastest route in is not networking, it is making things. Producers talk to producers
who have tracks. Even bad ones. Turn up somewhere with something you made and you have a
conversation; turn up with opinions and you have less of one.

\subsection{Realistic first steps}\label{li02_scene-realistic-first-steps}

\begin{enumerate}
\def\labelenumi{\arabic{enumi}.}
\tightlist
\item
  Go to ARC in September, or to Smartbar on a weeknight.
\item
  Go to Gramaphone and buy something on the staff\textquotesingle s recommendation.
\item
  Follow 5 Magazine\textquotesingle s calendar for a month to learn who the local promoters are.
\item
  Stand outside 206 S Jefferson St once, which is the Warehouse, landmarked in 2024. It is
  an ordinary building and that is somehow the point.
\end{enumerate}
